
The following is a specification for a subset of OCaml used during the
first part of the course.

\section{Literals}

\subsection*{Syntax}

\begin{align*}
  \bnfnt{int} &\bnfeq \bnft 0 \bnfalt \bnft 1 \bnfalt \bnft 2 \bnfalt \dots \\
  \bnfnt{float} &\bnfeq \bnft {0.} \bnfalt \bnft {1.} \bnfalt \bnft {2.} \bnfalt \dots \\
  \bnfnt{bool} &\bnfeq \bnft {true} \bnfalt \bnft {false} \\
  \bnfnt{expr} &\bnfeq \bnfnt{int} \bnfalt \bnfnt{float} \bnfalt \bnfnt{bool} \\
  \bnfnt{ty} &\bnfeq \bnft{int} \bnfalt \bnft{float} \bnfalt \bnft{bool}
\end{align*}
Literals are well-formed expressions by \textit{fiat}.  In the case of
Boolean literals, we can make the syntax explicit: there are two
literals, $\bnft{true}$ and $\bnft{false}$.  Infinite sets of literals
(e.g., in the case of integers and floating-point numbers) are defined
using regular expressions and are dealt with during \textit{lexical
  analysis}.  We additionally state that $\codel{int}$,
$\codel{float}$, and $\codel{bool}$ are well-formed types.

\begin{example}
  The integer literal $\codel{-2}$ and the Boolean literal
  $\codel{true}$ are well-formed expressions.
\end{example}

\subsection*{Typing}

\begin{mathpar}
  \inferrule*[right=intLit]{
    $\hl{$\codel{n}$ is an integer literal}$
  }{
    \G \ent \codel n : \intt
  }
  \and
  \inferrule*[right=floatLit]{
    $\hl{$\codel{n}$ is an float literal}$
  }{
    \G \ent \codel n : \floatt
  }
  \and
  \inferrule*[right=true]{
  }{
    \G \ent \truee : \boolt
  }
  \and
  \inferrule*[right=false]{
  }{
    \G \ent \falsee : \boolt
  }
\end{mathpar}
The rules for literals are \textit{axioms}, i.e., they have no
premises.  When building a derivation, these rules will appear as
leaves of the derivation tree.
\begin{itemize}
\item \textsc{intLit} says every integer literal is of type $\intt$ in any context.
\item \textsc{floatLit} says every floating-point literl is of type $\floatt$ in any context.
\item \textsc{true} says $\truee$ is of type $\boolt$ in any context.
\item \textsc{false} says $\falsee$ is of type $\boolt$ in any context.
\end{itemize}

\begin{example}
  The integer literal $\codel{35}$ is of type $\intt$ in the empty context, i.e.,
  the following is a valid derivation.
  \begin{prooftree}
    \AxiomC{}
    \RightLabel{\small \textsc{intLit}}
    \UnaryInfC{$\varnothing \ent \codel{35} : \intt$}
  \end{prooftree}
\end{example}

\subsection*{Semantics}

\begin{mathpar}
  \inferrule*[right=intLitE]{
    $\hl{$\codel{n}$ is an integer literal for $n$}$
  }{
    \codel{n} \evl n
  }
  \and
  \inferrule*[right=floatLitE]{
    $\hl{$\codel{n}$ is an float literal for $n$}$
  }{
    \codel n \evl n
  }
  \and
  \inferrule*[right=trueE]{
  }{
    \truee \evl \top
  }
  \and
  \inferrule*[right=falseE]{
  }{
    \falsee \evl \bot
  }
\end{mathpar}
\begin{itemize}
\item \textsc{intLitE} says an integer literal evaluates to the number it denotes.
\item \textsc{floatLitE} says a floating-point literal evaluates to the number it denotes.
\item \textsc{trueE} says the literal $\truee$ evaluates to the Boolean value $\top$.
\item \textsc{falseE} says the literal $\falsee$ evaluates to the Boolean value $\bot$.
\end{itemize}

\begin{example}
  The floating-point literal $\codel{3.14}$ evaluates to the number
  $3.14$, i.e., the following is a valid derivation.
  \begin{prooftree}
    \AxiomC{}
    \RightLabel{\small \textsc{floatLitE}}
    \UnaryInfC{$\codel{3.14} \evl 3.14$}
  \end{prooftree}
\end{example}

\section{Variables}

\subsection*{Syntax}

\begin{align*}
  \bnfnt{var} &\bnfeq \bnft{x} \bnfalt \bnft{foo\_bar} \bnfalt \bnft{abc'} \bnfalt \dots \\
  \bnfnt{expr} &\bnfeq \bnfnt{var}
\end{align*}
Variables is a well-formed expression.  As with number literals, we
can't make the syntax of a variable explicit, but we'll take a
variable to be alphanumeric with underscores and single-quotes, not
starting with a capital letter.

\subsection*{Typing}

\begin{mathpar}
  \inferrule*[right=var]{
    $\hl{$(x : \tau) \in \G$}$
  }{
    \G \ent x : \tau
  }
\end{mathpar}
\textsc{var} says if the variable $x$ is declared to be of type $\tau$
in the context $\G$, then $x$ is of type $\tau$ in the context $\G$.

\begin{example}
Since $(\codel{x} : \codel{int})$ appears in the context $\{\codel{x}
: \intt, \codel{y} : \boolt\}$, the variables $\codel{x}$ is of type $\intt$ in
this context, i.e., the following is a valid derivation.
\begin{prooftree}
  \AxiomC{}
  \RightLabel{\small \textsc{var}}
  \UnaryInfC{$\{ \codel x : \intt , \codel y : \boolt \} \ent \codel x : \intt$}
\end{prooftree}
\end{example}

\subsection*{Semantics}

There are no semantic rules for variables.  If evaluating an
expression requires evaluating a variable, then the expression is not
well-scoped.

\section{Operators}

\subsection*{Syntax}

\begin{align*}
  \bnfnt{op} &\bnfeq \bnft + \bnfalt \bnft - \bnfalt \bnft * \bnfalt \bnft / \bnfalt \bnft{mod} \\
  &\bnfaltn \bnft {+.} \bnfalt \bnft {-.} \bnfalt \bnft {*.} \bnfalt \bnft {/.} \\
  &\bnfaltn \bnft {\&\&} \bnfalt \bnft {||} \\
  &\bnfaltn \bnft {=} \bnfalt \bnft {<>} \bnfalt \bnft {<} \bnfalt \bnft {<=} \bnfalt \bnft {>} \bnfalt \bnft {>=} \\
  \bnfnt{expr} &\bnfeq \bnfnt{expr} \ \bnfnt{op} \ \bnfnt{expr}
\end{align*}
The first rule says the symbols $\codel{+}$, $\codel-$, $\codel *$,
$\codel/$, $\codel{mod}$, etc. are well-formed operators.  The second
says that if $\square$ is a well-formed operator and $e_1$ and $e_2$
are well-formed expressions, then $e_1 \ \square \ e_2$ is a
well-formed expression.

\subsection*{Typing}

\begin{mathpar}
  \inferrule*[right=addInt]{
    \G \ent e_1 : \intt
    \and
    \G \ent e_2 : \intt
  }{
    \G \ent e_1 \pluse e_2 : \intt
  }
  \and
  \inferrule*[right=subInt]{
    \G \ent e_1 : \intt
    \and
    \G \ent e_2 : \intt
  }{
    \G \ent e_1 \sube e_2 : \intt
  }
  \and
  \inferrule*[right=mulInt]{
    \G \ent e_1 : \intt
    \and
    \G \ent e_2 : \intt
  }{
    \G \ent e_1 \mule e_2 : \intt
  }
  \and
  \inferrule*[right=divInt]{
    \G \ent e_1 : \intt
    \and
    \G \ent e_2 : \intt
  }{
    \G \ent e_1 \dive e_2 : \intt
  }
  \and
  \inferrule*[right=modInt]{
    \G \ent e_1 : \intt
    \and
    \G \ent e_2 : \intt
  }{
    \G \ent e_1 \mode e_2 : \intt
  }
  \and
  \inferrule*[right=addFloat]{
    \G \ent e_1 : \floatt
    \and
    \G \ent e_2 : \floatt
  }{
    \G \ent e_1 \plusde e_2 : \floatt
  }
  \and
  \inferrule*[right=subFloat]{
    \G \ent e_1 : \floatt
    \and
    \G \ent e_2 : \floatt
  }{
    \G \ent e_1 \subde e_2 : \floatt
  }
  \and
  \inferrule*[right=mulFloat]{
    \G \ent e_1 : \floatt
    \and
    \G \ent e_2 : \floatt
  }{
    \G \ent e_1 \mulde e_2 : \floatt
  }
  \and
  \inferrule*[right=divFloat]{
    \G \ent e_1 : \floatt
    \and
    \G \ent e_2 : \floatt
  }{
    \G \ent e_1 \divde e_2 : \floatt
  }
  \and
  \inferrule*[right=and]{
    \G \ent e_1 : \boolt
    \and
    \G \ent e_2 : \boolt
  }{
    \G \ent e_1 \ande e_2 : \boolt
  }
  \and
  \inferrule*[right=or]{
    \G \ent e_1 : \boolt
    \and
    \G \ent e_2 : \boolt
  }{
    \G \ent e_1 \ore e_2 : \boolt
  }
  \and
  \inferrule*[right=eq]{
    \G \ent e_1 : \tau
    \and
    \G \ent e_2 : \tau
  }{
    \G \ent e_1 \eqe e_2 : \boolt
  }
  \and
  \inferrule*[right=neq]{
    \G \ent e_1 : \tau
    \and
    \G \ent e_2 : \tau
  }{
    \G \ent e_1 \neqe e_2 : \boolt
  }
  \and
  \inferrule*[right=lt]{
    \G \ent e_1 : \tau
    \and
    \G \ent e_2 : \tau
  }{
    \G \ent e_1 \lte e_2 : \boolt
  }
  \and
  \inferrule*[right=leq]{
    \G \ent e_1 : \tau
    \and
    \G \ent e_2 : \tau
  }{
    \G \ent e_1 \leqe e_2 : \boolt
  }
  \and
  \inferrule*[right=gt]{
    \G \ent e_1 : \tau
    \and
    \G \ent e_2 : \tau
  }{
    \G \ent e_1 \gte e_2 : \boolt
  }
  \and
  \inferrule*[right=geq]{
    \G \ent e_1 : \tau
    \and
    \G \ent e_2 : \tau
  }{
    \G \ent e_1 \geqe e_2 : \boolt
  }
\end{mathpar}
There are quite a few rules here, but in essence, they all say: if
$e_1$ and $e_2$ are some type in a given context then $e_1 \ \square
\ e_2$ is some type in the same context.

\begin{example}
  The following is derivation of the fact that $\codel{x} \pluse \codel{2}$ is of
  type $\intt$ in the context $\{ \codel x : \intt \}$.
  \begin{prooftree}
    \AxiomC{}
    \RightLabel{\small \textsc{var}}
    \UnaryInfC{$\{\codel x : \intt\} \ent \codel x : \intt$}
    \AxiomC{}
    \RightLabel{\small \textsc{intLit}}
    \UnaryInfC{$\{\codel x : \intt\} \ent \codel 2 : \intt$}
    \RightLabel{\small \textsc{addInt}}
    \BinaryInfC{$\{\codel x : \intt\} \ent \codel x \pluse \codel 2 : \intt$}
  \end{prooftree}
\end{example}

\subsection*{Semantics}

\begin{mathpar}
  \inferrule*[right=addIntE]{
    e_1 \evl v_1
    \and
    e_2 \evl v_2
    \and
    $\hl{$v = v_1 + v_2$}$
  }{
    e_1 \pluse e_2 \evl v
  }
  \and
  \inferrule*[right=subIntE]{
    e_1 \evl v_1
    \and
    e_2 \evl v_2
    \and
    $\hl{$v = v_1 - v_2$}$
  }{
    e_1 \sube e_2 \evl v
  }
  \and
  \inferrule*[right=mulIntE]{
    e_1 \evl v_1
    \and
    e_2 \evl v_2
    \and
    $\hl{$v = v_1 \times v_2$}$
  }{
    e_1 \mule e_2 \evl v
  }
  \and
  \inferrule*[right=divIntE]{
    e_1 \evl v_1
    \and
    e_2 \evl v_2
    \and
    $\hl{$v_2 \neq 0$}$
    \and
    $\hl{$v = v_1 / v_2$}$
  }{
    e_1 \dive e_2 \evl v
  }
  \and
  \inferrule*[right=modIntE]{
    e_1 \evl v_1
    \and
    e_2 \evl v_2
    \and
    $\hl{$v_2 \neq 0$}$
    \and
    $\hl{$v = v_1 \pmod {v_2}$}$
  }{
    e_1 \mode e_2 \evl v
  }
  \and
  \inferrule*[right=addFloatE]{
    e_1 \evl v_1
    \and
    e_2 \evl v_2
    \and
    $\hl{$v = v_1 + v_2$}$
  }{
    e_1 \pluse e_2 \evl v
  }
  \and
  \inferrule*[right=subFloatE]{
    e_1 \evl v_1
    \and
    e_2 \evl v_2
    \and
    $\hl{$v = v_1 - v_2$}$
  }{
    e_1 \sube e_2 \evl v
  }
  \and
  \inferrule*[right=mulFloatE]{
    e_1 \evl v_1
    \and
    e_2 \evl v_2
    \and
    $\hl{$v = v_1 \times v_2$}$
  }{
    e_1 \mule e_2 \evl v
  }
  \and
  \inferrule*[right=divFloatE]{
    e_1 \evl v_1
    \and
    e_2 \evl v_2
    \and
    $\hl{$v = v_1 / v_2$}$
  }{
    e_1 \dive e_2 \evl v
  }
  \and
  \inferrule*[right=andFalse]{
    e_1 \evl \bot
  }{
    e_1 \ande e_2 \evl \bot
  }
  \and
  \inferrule*[right=andTrue]{
    e_1 \evl \top
    \and
    e_2 \evl v
  }{
    e_1 \ande e_2 \evl v
  }
  \and
  \inferrule*[right=orTrue]{
    e_1 \evl \top
  }{
    e_1 \ore e_2 \evl \top
  }
  \and
  \inferrule*[right=orFalse]{
    e_1 \evl \bot
    \and
    e_2 \evl v
  }{
    e_1 \ore e_2 \evl v
  }
  \and
  \inferrule*[right=eqTrue]{
    e_1 \evl v_1
    \and
    e_2 \evl v_2
    \and
    $\hl{$v_1 = v_2$}$
  }{
    e_1 \eqe e_2 \evl \top
  }
  \and
  \inferrule*[right=eqFalse]{
    e_1 \evl v_1
    \and
    e_2 \evl v_2
    \and
    $\hl{$v_1 \neq v_2$}$
  }{
    e_1 \eqe e_2 \evl \bot
  }
  \and
  \inferrule*[right=neqTrue]{
    e_1 \evl v_1
    \and
    e_2 \evl v_2
    \and
    $\hl{$v_1 \neq v_2$}$
  }{
    e_1 \neqe e_2 \evl \top
  }
  \and
  \inferrule*[right=eqFalse]{
    e_1 \evl v_1
    \and
    e_2 \evl v_2
    \and
    $\hl{$v_1 = v_2$}$
  }{
    e_1 \neqe e_2 \evl \bot
  }
  \and
  \inferrule*[right=ltTrue]{
    e_1 \evl v_1
    \and
    e_2 \evl v_2
    \and
    $\hl{$v_1 < v_2$}$
  }{
    e_1 \lte e_2 \evl \top
  }
  \and
  \inferrule*[right=ltFalse]{
    e_1 \evl v_1
    \and
    e_2 \evl v_2
    \and
    $\hl{$v_1 \geq v_2$}$
  }{
    e_1 \lte e_2 \evl \bot
  }
  \and
  \inferrule*[right=lteTrue]{
    e_1 \evl v_1
    \and
    e_2 \evl v_2
    \and
    $\hl{$v_1 \leq v_2$}$
  }{
    e_1 \leqe e_2 \evl \top
  }
  \and
  \inferrule*[right=lteFalse]{
    e_1 \evl v_1
    \and
    e_2 \evl v_2
    \and
    $\hl{$v_1 > v_2$}$
  }{
    e_1 \leqe e_2 \evl \bot
  }
  \and
  \inferrule*[right=gtTrue]{
    e_1 \evl v_1
    \and
    e_2 \evl v_2
    \and
    $\hl{$v_1 > v_2$}$
  }{
    e_1 \gte e_2 \evl \top
  }
  \and
  \inferrule*[right=gtFalse]{
    e_1 \evl v_1
    \and
    e_2 \evl v_2
    \and
    $\hl{$v_1 \leq v_2$}$
  }{
    e_1 \gte e_2 \evl \bot
  }
  \and
  \inferrule*[right=geqTrue]{
    e_1 \evl v_1
    \and
    e_2 \evl v_2
    \and
    $\hl{$v_1 \geq v_2$}$
  }{
    e_1 \geqe e_2 \evl \top
  }
  \and
  \inferrule*[right=geqFalse]{
    e_1 \evl v_1
    \and
    e_2 \evl v_2
    \and
    $\hl{$v_1 < v_2$}$
  }{
    e_1 \geqe e_2 \evl \bot
  }
\end{mathpar}
Again, there are a lot of rules here, but they all in essence say that
if $e_1$ and $e_2$ evaluate to particular values, then $e_1 \ \square
\ e_2$ evaluates to the expected value. One case of note is Boolean
operators. Note how the rules enforce \textit{short-circuiting}.

\begin{example}
  The following is derivation of the fact that $\codel{2} \pluse
  \codel{3}$ evaluates to $5$.
  \begin{prooftree}
    \AxiomC{}
    \RightLabel{\small \textsc{intLitE}}
    \UnaryInfC{$\codel 2 \evl 2$}
    \AxiomC{}
    \RightLabel{\small \textsc{intLitE}}
    \UnaryInfC{$\codel 2 \evl 3$}
    \RightLabel{\small \textsc{addIntE}}
    \BinaryInfC{$\codel 2 \pluse \codel 3 \evl 5$}
  \end{prooftree}
\end{example}

\section{If-Expressions}

\subsection*{Syntax}

\begin{align*}
  \bnfnt{expr} \bnfeq \bnft{if} \ \bnfnt{expr} \ \bnft{then} \ \bnfnt{expr} \ \bnft{else} \ \bnfnt{expr}
\end{align*}
If $e_1$ is a well-formed expression and $e_2$ is a well-formed
expression and $e_3$ is a well-formed expression, then $\ife {e_1}
\thene {e_2} \elsee {e_3}$ is a well-formed expression.

\begin{example}
  Since $\codel 2$ is a well-formed expression, then $\ife \codel 2
  \thene \codel 2 \elsee \codel 2$ is a well-formed expression, i.e.,
  the following is a valid parse tree.
  \begin{center}
    \begin{forest}
      [$\bnfnt{expr}$
        [$\bnft{if}$]
        [$\bnfnt{expr}$
          [$\bnfnt{int}$
            [$\bnft{2}$]
          ]
        ]
        [$\bnft{then}$]
        [$\bnfnt{expr}$
          [$\bnfnt{int}$
            [$\bnft{2}$]
          ]
        ]
        [$\bnft{else}$]
        [$\bnfnt{expr}$
          [$\bnfnt{int}$
            [$\bnft{2}$]
          ]
        ]
      ]
    \end{forest}
  \end{center}
\end{example}

\subsection*{Typing}

\begin{mathpar}
  \inferrule*[right=if]{
    \G \ent e_1 : \boolt
    \and
    \G \ent e_2 : \tau
    \and
    \G \ent e_3 : \tau
  }{
    \G \ent \ife e_1 \thene e_2 \elsee e_3 : \tau
  }
\end{mathpar}
\textsc{if} says if $e_1$ is of type $\boolt$ in the context $\Gamma$
and $e_2$ is of type $\tau$ in the context $\Gamma$ and $e_3$ is
\textit{also} of type $\tau$ in the context $\Gamma$, then
$\ife{e_1} \thene {e_2} \elsee {e_3}$ is of type $\tau$ in the context $\Gamma$.

\begin{example}
  The following is a valid derivation of the fact that $\ife \falsee
  \thene \codel 2 \elsee \codel 3 \pluse \codel 4$ is of type $\intt$ in the empty
  context.

  \begin{prooftree}
    \AxiomC{}
    \RightLabel{\small \textsc{true}}
    \UnaryInfC{$\varnothing \ent \falsee : \boolt$}
    \AxiomC{}
    \RightLabel{\small \textsc{intLit}}
    \UnaryInfC{$\varnothing \ent \codel 2 : \intt$}
    \AxiomC{}
    \RightLabel{\small \textsc{intLit}}
    \UnaryInfC{$\varnothing \ent \codel 3 : \intt$}
    \AxiomC{}
    \RightLabel{\small \textsc{intLit}}
    \UnaryInfC{$\varnothing \ent \codel 4 : \intt$}
    \RightLabel{\small \textsc{addInt}}
    \BinaryInfC{$\varnothing \ent \codel 3 \pluse \codel 4 : \intt$}
    \RightLabel{\small \textsc{if}}
    \TrinaryInfC{$\varnothing \ent \ife \falsee \thene \codel 2 \elsee \codel 3 \pluse \codel 4 : \intt$}
  \end{prooftree}
\end{example}

\subsection*{Semantics}

\begin{mathpar}
  \inferrule*[right=ifTrue]{
    e_1 \evl \top
    \and
    e_2 \evl v
  }{
    \ife e_1 \thene e_2 \elsee e_3 \evl v
  }
  \and
  \inferrule*[right=ifFalse]{
    e_1 \evl \bot
    \and
    e_3 \evl v
  }{
    \ife e_1 \thene e_2 \elsee e_3 \evl v
  }
\end{mathpar}
\begin{itemize}
\item \textsc{ifTrue} says if $e_1$ evaluates to $\truee$ and $e_2$ evaluates to $v$, then $\ife e_1 \thene e_2 \elsee e_3$ evaluates to $v$
\item \textsc{ifFalse} says if $e_1$ evaluates to $\falsee$ and $e_3$ evaluates to $v$, then $\ife e_1 \thene e_2 \elsee e_3$ evaluates to $v$
\end{itemize}
In either case, we do not evaluate the irrelevant branch.

\begin{example}
  The following is a valid derivation of the fact that $\ife \falsee
  \thene \codel 2 \elsee \codel 3 \pluse \codel 4$ evaluates to $7$.
  \begin{prooftree}
    \AxiomC{}
    \RightLabel{\small \textsc{falseE}}
    \UnaryInfC{$\falsee \evl \bot$}
    \AxiomC{}
    \RightLabel{\small \textsc{intLitE}}
    \UnaryInfC{$\codel 3 \evl 3$}
    \AxiomC{}
    \RightLabel{\small \textsc{intLitE}}
    \UnaryInfC{$\codel 4 \evl 4$}
    \RightLabel{\small \textsc{addIntE}}
    \BinaryInfC{$\codel 3 \pluse \codel 4 \evl 7$}
    \RightLabel{\small \textsc{ifFalse}}
    \BinaryInfC{$\ife \falsee \thene \codel 2 \elsee \codel 3 \pluse \codel 4 \evl 7$}
  \end{prooftree}
\end{example}

\section{Let-Expressions}

\subsection*{Syntax}

\begin{align*}
  \bnfnt{expr} \bnfeq \bnft{let} \ \bnfnt{var} \ \bnft{=} \ \bnfnt{expr} \ \bnft{in} \ \bnfnt{expr}
\end{align*}
If $x$ is a variable and $e_1$ is a well-formed expression and $e_2$
is a well-formed expression, then $\lete x \eqe e_1 \ine e_2$ is a
well-formed expression.

\begin{example}
  Since $\codel{var\_name}$ is a well-formed variable name and $\codel
  2$ is a well-formed expression then $\lete \codel{var\_name} \eqe
  \codel 2 \ine \codel{var\_name}$ is a well-formed expression, i.e.,
  the following is a valid parse tree.

  \begin{center}
        \begin{forest}
      [$\bnfnt{expr}$
        [$\bnft{let}$]
        [$\bnfnt{var}$
          [$\bnft{var\_name}$]
        ]
        [$\bnft{=}$]
        [$\bnfnt{expr}$
          [$\bnfnt{int}$
            [$\bnft{2}$]
          ]
        ]
        [$\bnft{in}$]
        [$\bnfnt{expr}$
          [$\bnfnt{var}$
            [$\bnft{var\_name}$]
          ]
        ]
      ]
    \end{forest}
  \end{center}
\end{example}

\subsection*{Typing}

\begin{mathpar}
  \inferrule*[right=let]{
    \G \ent e_1 : \tau_1
    \and
    \G, x : \tau_1 \ent e_2 : \tau_2
  }{
    \G \ent \lete x \eqe e_1 \ine e_2 : \tau_2
  }
\end{mathpar}
\textsc{let} says if $e_1$ is of type $\tau_1$ in the context $\G$,
and $e_2$ is of type $\tau_2$ in the context $\G$ \textit{with the
  variable declaration $(x : \tau_1)$ added to it}, then
$\lete x \eqe e_1 \ine e_2$ is of type $\tau_2$ in the context $\G$.

\begin{example}
  The following is a valid derivation.

  \begin{prooftree}
    \AxiomC{}
    \RightLabel{\small \textsc{intLit}}
    \UnaryInfC{$\{\} \ent \codel 2 : \intt$}
    \AxiomC{}
    \RightLabel{\small \textsc{var}}
    \UnaryInfC{$\{\codel y : \intt\} \ent \codel y : \intt$}
    \AxiomC{}
    \RightLabel{\small \textsc{var}}
    \UnaryInfC{$\{\codel y : \intt\} \ent \codel y : \intt$}
    \RightLabel{\small \textsc{intAdd}}
    \BinaryInfC{$\{\codel y : \intt\} \ent \codel y \pluse \codel y : \intt$}
    \RightLabel{\small \textsc{let}}
    \BinaryInfC{$\{\} \ent \lete \codel y \eqe \codel 2 \ine \codel y \pluse \codel y : \intt$}
  \end{prooftree}
\end{example}

\subsection*{Semantics}

\begin{mathpar}
  \inferrule*[right=letE]{
    e_1 \evl v_1
    \and
    $\hl{$[v_1 / x]e_2 = e'$}$
    \and
    e' \evl v
  }{
    \lete x \eqe e_1 \ine e_2 \evl v
  }
\end{mathpar}
\textsc{letE} says if $e_1$ evaluates to $v_1$ and $e_2$ \textit{with
  $v_1$ substituted for $x$} evaluates to $v$, then $\lete x \eqe e_1
\ine e_2$ evaluates to $v$.

\begin{example}
  The following is a valid derivation.
  \begin{prooftree}
    \AxiomC{}
    \RightLabel{\small \textsc{intLitE}}
    \UnaryInfC{$\codel 2 \evl 2$}
    \AxiomC{}
    \RightLabel{\small \textsc{intLitE}}
    \UnaryInfC{$\codel 2 \evl 2$}
    \AxiomC{}
    \RightLabel{\small \textsc{intLitE}}
    \UnaryInfC{$\codel 2 \evl 2$}
    \RightLabel{\small \textsc{intAddE}}
    \BinaryInfC{$\codel 2 \pluse \codel 2 \evl 4$}
    \RightLabel{\small \textsc{letE}}
    \BinaryInfC{$\lete \codel y \eqe \codel 2 \ine \codel y \pluse \codel y \evl 4$}
  \end{prooftree}
\end{example}

\section{Functions}

\subsection*{Syntax}

\begin{align*}
  \bnfnt{expr} \bnfeq \bnft{fun} \ \bnfnt{var} \ \bnft{->} \ \bnfnt{expr}
\end{align*}
If $x$ is a well-formed variable and $e$ is a well-formed expression,
then $\fune x \toe e$ is a well-formed expression.

\begin{example}
  The following is a valid parse tree for the well-formed expression
  $\fune \codel{name} \toe \codel{name} \pluse \codel 2$.
  \begin{center}
    \begin{forest}
      [$\bnfnt{expr}$
        [$\bnft{fun}$]
        [$\bnfnt{var}$
          [$\bnft{name}$]
        ]
        [$\bnft{->}$]
        [$\bnfnt{expr}$
          [$\bnfnt{var}$
            [$\bnft{name}$]
          ]
          [$\bnft{+}$]
          [$\bnfnt{int}$
            [$\bnft{3}$]
          ]
        ]
      ]
    \end{forest}
  \end{center}
\end{example}

\subsection*{Typing}

\begin{mathpar}
  \inferrule*[right=fun]{
    \G, x : \tau_1 \ent e : \tau_2
  }{
    \G \ent \fune x \toe e : \tau_1 \toe \tau_2
  }
\end{mathpar}
\textsc{fun} says if $e$ is of type $\tau_2$ in the context $\Gamma$
\textit{together with the variable declaration $(x : \tau_1)$ added to
  it}, then $\fune {x} \toe {e}$ is of type ${\tau_1} \toe {\tau_2}$.

\begin{example}
  The following is a valid derivation.
  \begin{prooftree}
    \AxiomC{}
    \RightLabel{\small \textsc{intLit}}
    \UnaryInfC{$\{ \codel z : \floatt \} \ent \codel 2 : \intt$}
    \RightLabel{\small \textsc{fun}}
    \UnaryInfC{$\{\} \ent \fune \codel z \toe \codel 2 : \floatt \toe \intt$}
  \end{prooftree}
\end{example}

\subsection*{Semantics}

\begin{mathpar}
  \inferrule*[right=funE]{
  }{
    \fune x \toe e \evl \lambda x . e
  }
\end{mathpar}
\textsc{funE} says a function $\fune x \toe e$ evaluates to a special
function value that essentially hold the variables $x$ and the body
$e$ with the notation $\lambda x . e$.

\begin{example}
  The following is a valid derivation.
  \begin{prooftree}
    \AxiomC{}
    \RightLabel{\small \textsc{funEval}}
    \UnaryInfC{$\fune \codel x \toe \codel x \pluse \codel x \evl \lambda \codel x . \codel x \pluse \codel x$}
  \end{prooftree}
\end{example}

\section{Application}

\subsection*{Syntax}

\begin{align*}
  \bnfnt{expr} \bnfeq \bnfnt{expr} \ \bnfnt{expr}
\end{align*}
If $e_1$ is a well-formed expression and $e_2$ is a well-formed
expression, then $e_1 \ e_2$ is a well-formed expression.

\subsection*{Typing}

\begin{mathpar}
  \inferrule*[right=app]{
    \G \ent e_1 : \tau_2 \toe \tau
    \and
    \G \ent e_2 : \tau_2
  }{
    \G \ent e_1 \ e_2 : \tau
  }
\end{mathpar}
\textsc{app} says if $e_1$ is of type $\tau_2 \toe \tau$ and $e_2$ is of type $\tau_2$,
then $e_1 \ e_2$ is of type $\tau$.

\begin{example}
  The following is a valid derivation.
  \begin{prooftree}
    \AxiomC{}
    \RightLabel{\small \textsc{intLit}}
    \UnaryInfC{$\{ \codel z : \floatt \} \ent \codel 2 : \intt$}
    \RightLabel{\small \textsc{fun}}
    \UnaryInfC{$\{\} \ent \fune {\codel z} \toe {\codel 2} : \floatt \toe \intt$}
    \AxiomC{}
    \RightLabel{\small \textsc{floatLit}}
    \UnaryInfC{$\{\}\ent \codel{2.0} : \floatt$}
    \RightLabel{\small \textsc{app}}
    \BinaryInfC{$\{\} \ent (\fune \codel z \toe \codel 2) \ \codel{2.0} : \intt$}
  \end{prooftree}
\end{example}

\subsection*{Semantics}

\begin{mathpar}
  \inferrule*[right=appE]{
    e_1 \evl \lambda x . e
    \and
    e_2 \evl v_2
    \and
    $\hl{$[v_2 / x]e = e'$}$
    \and
    e' \evl v
  }{
    e_1 \ e_2 \evl v
  }
\end{mathpar}
\textsc{appE} says if $e_1$ evaluates to the function value $\lambda x
. e$ and $e_2$ evaluates to $v_2$ and $e$ \textit{with $v_2$ substituted for
$x$} evaluates to $v$, then $e_1 \ e_2$ evaluates to $v$.

\begin{example}
  The following is a valid derivation.
  \begin{prooftree}
    \AxiomC{}
    \RightLabel{\small \textsc{funEval}}
    \UnaryInfC{$\fune \codel x \toe \codel x \pluse \codel x \evl \lambda \codel x . \codel x \pluse \codel x$}
    \AxiomC{}
    \RightLabel{\small \textsc{intLitE}}
    \UnaryInfC{$\codel 2 \evl 2$}
    \AxiomC{}
    \RightLabel{\small \textsc{intLitE}}
    \UnaryInfC{$\codel 2 \evl 2$}
    \AxiomC{}
    \RightLabel{\small \textsc{intLitE}}
    \UnaryInfC{$\codel 2 \evl 2$}
    \RightLabel{\small \textsc{addIntE}}
    \BinaryInfC{$\codel 2 \pluse \codel 2 \evl 4$}
    \RightLabel{\small \textsc{appE}}
    \TrinaryInfC{$(\fune \codel x \toe \codel x \pluse \codel x) \ \codel 2 \evl 4$}
  \end{prooftree}
\end{example}

\begin{remark}
  (\textbf{IMPORTANT}) There is a subtlety to substitution: when a
  value is substituted for a variable in an expression, \textit{it is
    converted back into an expression}, e.g., $[2 / \codel x](\codel x
  \pluse \codel x) = \codel 2 \pluse \codel 2$.  This will be made
  more clear when we see how this manifests in our implementation of
  interpreters.  For now, it is important to note that it is simple to
  convert the values we care about back into expressions via the
  following mapping:
  \begin{align*}
    \textit{integer} &\mapsto \codel{integer literal} & \text{(e.g., $2 \mapsto \codel{2}$)} \\
    \textit{float} &\mapsto \codel{float literal} & \text{(e.g., $3.5 \mapsto \codel{3.5}$)} \\
    \top &\mapsto \truee \\
    \bot &\mapsto \falsee \\
    \lambda x . e & \mapsto \fune x \toe e
  \end{align*}
\end{remark}
